\chapter{อุณหพลศาสตร์และทฤษฎีจลน์ของก๊าซ}
\label{ch:ch10_thermodynamics_gases}
\index{อุณหพลศาสตร์และก๊าซ}

\begin{conceptbox}[title=วัตถุประสงค์การเรียนรู้ประจำบท]
เมื่อสิ้นสุดการศึกษาบทเรียนนี้ นักศึกษาสามารถ
\begin{enumerate}
    \item เมื่อศึกษาบทนี้จบแล้ว ผู้เรียนควรมีความสามารถดังนี้
    \item อธิบายกฎของแก๊สอุดมคติ (Ideal Gas Law) และแบบจำลองจลน์ของแก๊ส (Kinetic Theory of Gases) ได้อย่างถูกต้อง
    \item อธิบายความหมายของพลังงานภายใน (Internal Energy) ความร้อน (Heat) และงาน (Work) ในบริบทของระบบทางอุณหพลศาสตร์
    \item อธิบายและประยุกต์ใช้กฎข้อที่หนึ่งของอุณหพลศาสตร์ในการวิเคราะห์กระบวนการต่างๆ
    \item วิเคราะห์และคำนวณงาน ความร้อน และการเปลี่ยนแปลงพลังงานภายในของแก๊สในกระบวนการ isothermal, isobaric, isochoric และ adiabatic
\end{enumerate}
\end{conceptbox}

\section{ทฤษฎีและหลักการพื้นฐาน}

\begin{bioappbox}
\textbf{เอนโทรปี ชีวฟิสิกส์ของเยื่อหุ้มเซลล์ และการแลกเปลี่ยนก๊าซ}

กฎข้อที่หนึ่งของอุณหพลศาสตร์ $\Delta U = Q - W$ อธิบายงบประมาณพลังงานในสิ่งมีชีวิต ขณะที่กฎข้อที่สองว่าด้วยเอนโทรปี ($\Delta S \ge 0$) อธิบายว่าสิ่งมีชีวิตต้องบริโภคพลังงานอิสระ (Gibbs Free Energy, $\Delta G = \Delta H - T\Delta S$) เพื่อรักษาสภาพระเบียบของระบบเซลล์มิให้พังทลายลง นอกจากนี้ กฎของก๊าซในอุดมคติ $PV = nRT$ และกฎของดอลตันเรื่องความดันย่อย ใช้อธิบายกลไกการแพร่ของก๊าซออกซิเจนและคาร์บอนไดออกไซด์ผ่านถุงลมปอดและปากใบของพืช
\end{bioappbox}

\begin{labbox}
1. ตรวจสอบตำแหน่งศูนย์ (Zero Error) ของเครื่องมือวัดทุกชนิดก่อนเริ่มบันทึกข้อมูล
2. ทำการวัดซ้ำอย่างน้อย 3 ถึง 5 ครั้งในแต่ละจุดทดลองเพื่อคำนวณหาค่าเฉลี่ยและค่าเบี่ยงเบนมาตรฐาน
3. บันทึกผลการวัดพร้อมระบุหน่วยเอสไอ (SI Units) และค่าความคลาดเคลื่อนของการวัดเสมอ
\end{labbox}

\section{ตัวอย่างโจทย์และการวิเคราะห์เชิงฟิสิกส์}

\begin{examplebox}
ตัวอย่างที่ 10.1

แก๊สอุดมคติจำนวน 2.0 โมล บรรจุอยู่ในภาชนะที่มีปริมาตร 0.050 m$^3$ ที่อุณหภูมิ 300 K จงหาความดันของแก๊สนี้

วิธีทำ

จากกฎของแก๊สอุดมคติ PV = nRT
\end{examplebox}

\section{แบบฝึกหัดทบทวนและคำถามท้ายบท}

\begin{enumerate}
    \item จงอธิบายสมมติฐานพื้นฐานของแบบจำลองแก๊สอุดมคติมาอย่างน้อย 3 ข้อ
    \item แก๊สอุดมคติจำนวน 3.0 โมล บรรจุในภาชนะปริมาตร 0.10 m$^3$ ที่อุณหภูมิ 350 K จงหาความดันของแก๊ส
    \item แก๊สในภาชนะหนึ่งมีความดัน 1.5 $\times$ 10⁵ Pa ปริมาตร 0.020 m$^3$ และมีจำนวน 1.2 โมล จงหาอุณหภูมิของแก๊ส
    \item จงอธิบายความแตกต่างระหว่างความร้อน (Q) พลังงานภายใน (U) และงาน (W) ในบริบทของกฎข้อที่หนึ่งของอุณหพลศาสตร์
    \item แก๊สขยายตัวภายใต้ความดันคงที่ 1.0 $\times$ 10⁵ Pa จากปริมาตร 0.0020 m$^3$ เป็น 0.0060 m$^3$ จงหางานที่แก๊สกระทำ
    \item ระบบหนึ่งได้รับความร้อน 400 J และมีงานกระทำต่อระบบ (โดยสิ่งแวดล้อม) เท่ากับ 150 J จงหาการเปลี่ยนแปลงพลังงานภายในของระบบ (คำนวณโดยแปลงงานที่กระทำต่อระบบเป็น W ในสมการ $\Delta$U = Q − W ก่อน)
\end{enumerate}

% สิ้นสุดบทเรียน
